\documentclass[12pt]{article}

\usepackage{sbc-template}
\usepackage{graphicx,url}
\usepackage[utf8]{inputenc}
\usepackage{amsmath}

\sloppy

\title{Mining Biosynthetic Gene Clusters in \textit{Helicobacter pylori} Genomes Associated with Premalignant and Malignant Gastric Phenotypes: Preliminary Results and Next Steps}

\author{Carlos Eugênio Costa de Lima, Wilson Araújo da Silva Junior}

\address{Laboratory of Molecular Genetics and Omics (LGMO)\\
Faculty of Medicine of Ribeirão Preto -- University of São Paulo (USP)\\
Ribeirão Preto -- SP -- Brazil
\email{carlos.eugenio.cl@usp.br, wilsonjr@usp.br}
}

\begin{document}

\maketitle

\begin{abstract}
Biosynthetic gene clusters (BGCs) encode microbial secondary metabolites, several of which have documented roles in tumor biology, yet BGC mining has never been applied to \textit{Helicobacter pylori} in the context of gastric carcinogenesis. Building on an unpublished undergraduate pilot study by the same group, which mined BGCs from two metagenome-assembled \textit{H. pylori} genomes, this work-in-progress applies the same funcscan-based strategy (antiSMASH, DeepBGC, GECCO) to 554 preassembled \textit{H. pylori} genomes from cancer and precancerous cohorts, followed by domain-based Gene Cluster Family (GCF) clustering and statistical testing (Fisher's exact test, CMH, and logistic regression, corrected for multiple testing across the full tested set). Preliminary results show a subtle but significant increase in BGC richness in cancer genomes, with no individual GCF reaching significance after correction, motivating the next steps described here.
\end{abstract}

\section{Introduction}

Biosynthetic gene clusters (BGCs) are co-located sets of genes responsible for the production of microbial secondary metabolites. Genome-mining tools such as antiSMASH and DeepBGC enable their identification and classification from genomic and metagenomic data \cite{blin2025,hannigan2019}. Although BGCs have already been linked to antitumor agents in other bacterial genera, no study to date has directly investigated bacterial BGCs in the context of the gastric cancer microbiome; the existing literature on this topic focuses on broad-spectrum taxonomic and functional profiling, without incorporating the biosynthetic perspective \cite{coker2018,noto2017}.

In an unpublished undergraduate thesis by this group\footnote{Lima, C.\ E.\ C.\ \textit{Identificação de Biosynthetic Gene Clusters (BGCs) associados à carcinogênese em histopatologias gástricas}. Undergraduate thesis, Faculty of Medicine of Ribeirão Preto, University of São Paulo, 2025 (unpublished).}, BGC mining was applied, as a pilot, to two metagenome-assembled \textit{H. pylori} genomes (MAGs) recovered from gastric biopsies of two high-risk gastric cancer patients. Despite the very low fraction of recoverable microbial DNA in human samples (about 2.2\% of filtered reads), 28 biosynthetic regions of interest were identified, including modules linked to the \textit{de novo} lipogenesis pathway (acetyl-CoA carboxylase, beta-ketoacyl-ACP synthase, 3-oxoacyl-ACP reductase) and genes involved in tRNA modification and capsular polysaccharide biosynthesis, several of which have documented parallels in the tumor progression literature. This pilot result motivated the present stage of the project, a Master's thesis, which forgoes \textit{de novo} metagenomic assembly -- heavily limited at scale given the low bacterial coverage depth achievable in human tissue -- in favor of mining BGCs from already-sequenced, publicly deposited \textit{H. pylori} genomes, enabling a formal statistical comparison between clinical phenotypes with substantially greater sample size.

\section{Methods}

We obtained 554 preassembled \textit{H. pylori} genomes from NCBI, 179 from gastric-cancer-related samples and 375 from precursor cases (chronic gastritis and intestinal metaplasia). BGC mining was performed with the \texttt{nf-core/funcscan} pipeline (antiSMASH 8.0.1, DeepBGC 0.1.31, and GECCO 0.9.10), with results aggregated by comBGC into a single, harmonized table of Pfam domains per BGC \cite{ewels2020,ditommaso2017,frangenberg2024}.

Since BiG-SCAPE requires antiSMASH-formatted GenBank files -- incompatible with DeepBGC and GECCO outputs -- we implemented, using Python and Bash scripts, a tool-agnostic Gene Cluster Family (GCF) clustering workflow operating directly on the harmonized Pfam domain arrays. Prior to similarity computation, DeepBGC BGCs with probability $<$ 0.50 were excluded, and BGCs without an assigned product class (unclassified DeepBGC predictions and GECCO ``Unknown'') were retained as a separate class. Pairwise similarity followed an adaptation of \cite{robey2021} under two regimes: for BGCs with identifiable backbone domains,
$$\text{Score} = \sqrt{0.8 \times \text{sequence identity (hmmalign)} + 0.2 \times \text{Jaccard}};$$
for the remaining BGCs, $\text{Score} = \sqrt{0.2 \times \text{Jaccard}}$. Final GCF assignment used DBSCAN ($\varepsilon = 0.50$, \texttt{min\_samples} = 2) on the resulting distance matrix (1 -- Score).

For statistical testing, every non-singleton GCF carried by at least 3 genomes (a criterion set \textit{a priori}) was tested for association with clinical phenotype (cancer vs. non-cancer) using Fisher's exact test, the Cochran--Mantel--Haenszel (CMH) test stratified by continent, and multivariable logistic regression (adjusted for continent, age, and total BGC richness per genome), with Benjamini--Hochberg (FDR) correction applied across the entire tested set \cite{benjaminihochberg1995}.

\section{Preliminary Results}

Funcscan mining across the 554 genomes yielded 2,078 BGC predictions before filtering (2,076 after), of which 1,518 (73.1\%) were grouped into 92 non-singleton GCFs and 558 remained singletons. The largest GCF (GCF\_0001, 1,003 BGCs) corresponds to the universal terpene-precursor locus detected by antiSMASH in virtually every genome; the median size of the remaining GCFs was only 3 BGCs, reflecting the atypical and restricted biosynthetic repertoire of \textit{H. pylori} already documented in the pilot study.

BGC richness per genome was modestly, but significantly, higher in cancer genomes overall (median of 4 BGCs/genome vs. 3 in non-cancer; Wilcoxon $W$ = 40459.5; $p$ = 5.96$\times10^{-5}$) -- a subtle difference in absolute terms, but a statistically robust one. Breaking this down by prediction tool: antiSMASH showed virtually identical distributions between Cancer and Non-cancer ($W$ = 33898.5; $p_{adj}$ = 0.387; ns), consistent with this tool typically detecting only 1--2 BGCs per genome; DeepBGC showed a significant difference ($W$ = 37563.0; $p$ = 1.30$\times10^{-5}$; $p_{adj}$ = 2.61$\times10^{-5}$; ***); GECCO could not be evaluated due to its low prediction frequency.

GCF-level richness per genome also differed significantly by the Wilcoxon test ($W$ = 38477; $p$ = 0.00257), even though the median was identical (2 GCFs/genome) in both groups. We flag this result with a caveat: because the Wilcoxon test is sensitive to distributional shape beyond the median (dispersion, ties, skew), a significant $p$-value alongside equal medians should not, on its own, be read as evidence of a clinically meaningful shift in typical GCF richness. We report it here as a distributional signal that warrants further characterization -- for instance via an effect-size estimate or a closer look at the tails of the distribution -- rather than as a straightforward ``cancer genomes carry more GCFs'' finding.

At the individual-GCF level, testing the 49 GCFs (out of 91 non-singleton GCFs in total) that met the pre-specified $\geq$ 3-carrier-genome criterion, Fisher's exact test with Benjamini--Hochberg correction across the full tested set found no GCF with $p_{adj} < 0.05$ or even $p_{adj} < 0.10$ (smallest observed $p_{adj}$: 0.109). Of these 49 eligible GCFs, 36 yielded a computable continent-stratified CMH statistic and all 49 yielded a computable multivariable logistic regression estimate; none reached adjusted significance by either method (0/36 and 0/49 respectively, $p_{adj} < 0.05$). In other words, despite the aggregate richness signal described above, no individual GCF carries a statistically robust association with cancer phenotype in this cohort once the number of GCFs actually tested is properly accounted for.

A quasi-Poisson model (BGC richness $\sim$ phenotype + age) showed additive, independent effects: IRR = 1.119 for cancer (95\% CI: 1.000--1.250; $p$ = 0.0487) and IRR = 1.005 per year of age (95\% CI: 1.001--1.008; $p$ = 0.0047), with no evidence of a phenotype $\times$ age interaction (likelihood-ratio test, $p$ = 0.864). Spearman correlations between age and BGC richness were weak but significant overall ($\rho$ = 0.16; $p$ = 1.14$\times10^{-4}$) and in the non-cancer subgroup ($\rho$ = 0.15; $p$ = 0.0048), but not significant in the cancer subgroup alone ($\rho$ = 0.08; $p$ = 0.305); this does not contradict the absence of an interaction detected by the pooled model.

\section{Discussion}

Taken together, the preliminary results point to a subtle but statistically robust increase in overall BGC/GCF richness in cancer genomes, while no individual GCF carries a significant association after appropriate correction for the number of GCFs tested. This is consistent with the hypothesis, already raised in the pilot study, that \textit{H. pylori} carries a restricted and highly conserved biosynthetic repertoire, in which phenotype-related effects may manifest more as quantitative differences in overall biosynthetic load than as the presence or absence of specific families. The breakdown by prediction tool (antiSMASH ns vs. DeepBGC significant) is reported here for methodological transparency, and it reinforces the value of combining rule-based and deep-learning-based mining approaches when working with an organism whose secondary metabolism is atypical \cite{zhu2025}.

\section{Next Steps}

Planned next steps for this Master's project include: (i) completing \texttt{hmmalign}-based backbone-domain identification for all eligible BGC pairs, currently pending due to a computational-environment issue (\texttt{hmmpress}/Pfam-A.hmm); (ii) an epsilon sensitivity sweep for DBSCAN, as a diagnostic prior to a full investment in the \texttt{hmmalign}-based solution, to assess whether the current GCF granularity is overly permissive or reflects genuine within-family expansion; (iii) correlating \textit{H. pylori} phylogeny with the observed GCF repertoire; (iv) searching for a malignancy signature based on combinations of GCFs/BGCs, beyond the univariate GCF-by-GCF testing already performed; and (v) discussing an experimental approach for functional validation of the most relevant BGCs in an animal model.

\section{Conclusion}

This work-in-progress extends, to a cohort of 554 public \textit{H. pylori} genomes, the BGC mining approach validated as a pilot on two metagenome-assembled genomes. Preliminary results indicate a subtle but statistically robust signal of higher biosynthetic richness in genomes associated with gastric cancer, and no individual GCF association surviving correction for multiple testing across the full eligible set. These findings, even though negative at the individual-GCF level, consolidate the clustering and pre-specified statistical testing methodology as a solid foundation for the phylogeny, combined malignancy signature, and functional validation analyses planned for the completion of the project.

\bibliographystyle{sbc}
\bibliography{references}

\end{document}
